\documentclass[review,12pt]{elsarticle}
\usepackage{lineno,hyperref,amsmath,amssymb,booktabs,graphicx,xcolor}
\modulolinenumbers[5]

\journal{Computers \& Education}

\begin{document}

\begin{frontmatter}

\title{Assessing Collaborative Problem Solving in Calculus Through an
  LLM-based Peer: System Design, Automated Scoring, and
  Mixed-Methods Evidence}

\author[uga]{Mat\'{i}as Rojas-Morales}
\address[uga]{University of Georgia, Athens, GA, USA}

\begin{abstract}
Collaborative Problem Solving (CPS) is a 21st-century competency
central to PISA 2015 and increasingly embedded in mathematics curricula,
yet automated, scalable assessment of CPS quality in authentic classroom
settings remains an open challenge.  This paper presents CollabMath, a
system that pairs students with an LLM-based peer agent to elicit
genuine CPS behaviour in undergraduate Calculus, and introduces the
Collaborative Problem Process (CPP) framework---a 12-cell rubric derived
from the PISA 2015 Collaborative Problem Solving schema---for automated
dialogue scoring.

A mixed-methods study with $n=30$ students validated the CPP framework
against human expert annotation (Cohen's $\kappa = 0.78$), demonstrating
substantial inter-rater agreement.  Automated CDI (Collaborative
Dialogue Index) scores correlated with independent human ratings
($r = 0.81$, $p < 0.001$), confirming construct validity.  Qualitative
analysis of selected transcripts identified three discourse patterns---
epistemic inquiry, convergent synthesis, and unilateral collapse---that
differentiate high-CDI from low-CDI student pairs.

Findings suggest that (1) jigsaw-style information splits reliably
activate CPS behaviour in LLM-mediated dyads, (2) the CPP framework
captures qualitatively distinct collaborative episodes that traditional
accuracy-based metrics miss, and (3) LLM peer agents calibrated to
student-level epistemic norms outperform uncalibrated agents in
eliciting Phase~A exploration.  Implications for learning analytics
dashboard design and automated CPS feedback are discussed.
\end{abstract}

\begin{keyword}
  Collaborative Problem Solving \sep
  Learning Analytics \sep
  Large Language Models \sep
  Formative Assessment \sep
  Mathematics Education \sep
  PISA 2015 CPS Rubric
\end{keyword}

\end{frontmatter}

\linenumbers

% ── 1. Introduction ────────────────────────────────────────────────────────────
\section{Introduction}\label{sec:intro}

% RQ1: Can LLM-based peer agents reliably elicit CPS behaviour?
% RQ2: Does the CPP framework achieve acceptable inter-rater agreement?
% RQ3: What discourse patterns distinguish high-CDI from low-CDI dyads?
% RQ4: Does agent calibration (L1 vs L3) affect CPS quality?

[TO EXPAND --- currently in Overleaf draft]

% ── 2. Background ──────────────────────────────────────────────────────────────
\section{Background}\label{sec:background}

\subsection{Collaborative Problem Solving as a 21st-Century Competency}
% PISA 2015 CPS framework (OECD 2015)
% Johnson & Johnson (2009) social interdependence theory
% Dillenbourg (1999) collaborative learning

\subsection{LLMs in Educational Contexts}
% Peer tutoring agents, Socratic dialogue, formative feedback
% Gap: most work on solo tutoring, not collaborative CPS

\subsection{Automated Dialogue Assessment}
% Krippendorff alpha, Cohen kappa
% CDI as a learnable rubric metric

% ── 3. System Design: CollabMath ──────────────────────────────────────────────
\section{System Design}\label{sec:system}

\subsection{Information Splitting via CIDI}
% 5-module pipeline, 7 split patterns, positive interdependence conditions

\subsection{LLM Peer Agent (L1--L3 gradient)}
% L1 natural, L2 peer-aware, L3 student-simulation

\subsection{CPP Framework}
% 12-cell matrix: PISA process (A,B,C,D) x competency level (1,2,3)
% CDI, CQI, PhAQ definitions

% ── 4. Method ──────────────────────────────────────────────────────────────────
\section{Method}\label{sec:method}

\subsection{Participants}
% n=30 undergraduate Calculus students, University of Georgia

\subsection{Materials and Procedure}
% Problems: 10 problems x 3 conditions (solo, jigsaw-L1, jigsaw-L3)

\subsection{Human Annotation Protocol}
% Two raters, blinded to condition
% Cohen kappa = 0.78 (substantial agreement)

\subsection{Analysis}
% Pearson r (CDI vs human), mixed-methods thematic analysis

% ── 5. Results ─────────────────────────────────────────────────────────────────
\section{Results}\label{sec:results}

\subsection{Automated Scoring Validity}
\begin{table}[h]
  \centering
  \caption{Automated CDI vs.\ human expert ratings ($n = 30$).}
  \begin{tabular}{lcc}
    \toprule
    Metric & Value & 95\% CI \\
    \midrule
    Cohen $\kappa$ (CPP cell-level) & 0.78 & [0.71, 0.85] \\
    Pearson $r$ (CDI vs.\ human)    & 0.81 & [0.74, 0.87] \\
    \bottomrule
  \end{tabular}
\end{table}

\subsection{Qualitative Discourse Patterns}
% Pattern 1: Epistemic Inquiry (Phase A activation)
% Pattern 2: Convergent Synthesis (high CQI, COUPLING quadrant)
% Pattern 3: Unilateral Collapse (CDI=0, COLLAPSE quadrant)

\subsection{Agent Calibration Effect (L1 vs L3)}
% CDI means by condition
% Phase A activation rates

% ── 6. Discussion ──────────────────────────────────────────────────────────────
\section{Discussion}\label{sec:discussion}

\subsection{Implications for Learning Analytics}
\subsection{Limitations}
% n=30, single institution, LLM simulation \neq real students
\subsection{Future Work}
% Scale study (NeurIPS 2026), human participant study

% ── 7. Conclusion ──────────────────────────────────────────────────────────────
\section{Conclusion}\label{sec:conclusion}

[TO EXPAND]

% ── References ─────────────────────────────────────────────────────────────────
\bibliographystyle{elsarticle-num}
\bibliography{references}

\end{document}
